% Document template suitable for use as a LaTeX master-file 
% for thesis works in University of Turku Department of Computing
%
% Technical usage guide: https://tech.utugit.fi/soft/thesis/doc/doc/overview/
% 

\documentclass[language=finnish,version=final,mainfont=none,sharelatex=false]{utuftthesis}
\setcounter{secnumdepth}{2}
\setcounter{tocdepth}{2}
\usepackage{float}
\usepackage[caption=false]{subfig}
\usepackage{graphicx}  % For \resizebox
\usepackage{array}     % For adjusting column widths
\usepackage{lscape}    % For landscape environment

% Define the algorithm environment
%\makeatletter
\providecommand\textquotedblplain{%
  \bgroup\addfontfeatures{Mapping=}\char34\egroup}
\providecommand{\tabularnewline}{\\}
\floatstyle{ruled}
\newfloat{algorithm}{tbp}{loa}
\providecommand{\algorithmname}{Algoritmi}
\floatname{algorithm}{\protect\algorithmname}
%\makeatother

\addbibresource{Bibliografia.bib}

\begin{document}

\pubyear{2023}
\pubmonth{11}
\publab{Tietotekniikka}
\publaben{Information Technology}
\pubtype{tkk}
\title{Raportointityökalun kehittäminen Sovelia® PLM-järjestelmään}
\author{Elias Peltonen}

\maketitle

\keywords{tähän, lista, avainsanoista}
% TODO: good/bad keywords

\keywordsen{here, a, list, of, keywords}
\begin{abstract}
Tarkempia ohjeita tiivistelmäsivun laadintaan läytyy opiskelijan yleisoppaasta,
josta alla lyhyt katkelma.

Bibliografisten tietojen jälkeen kirjoitetaan varsinainen tiivistelmä.
Sen on oletettava, että lukijalla on yleiset tiedot aiheesta. Tiivistelmän
tulee olla ymmärrettävissä ilman tarvetta perehtyä koko tutkielmaan.
Se on kirjoitettava täydellisinä virkkeinä, väliotsakeluettelona.
On käytettävä vakiintuneita termejä. Viittauksia ja lainauksia tiivistelmään
ei saa sisällyttää, eikä myäskään tietoja tai väitteitä, jotka eivät
sisälly itse tutkimukseen. Tiivistelmän on oltava mahdollisimman ytimekäs
n. 120–250 sanan pituinen itsenäinen kokonaisuus, joka mahtuu ykkäsvälillä
kirjoitettuna vaivatta yhdelle tiivistelmäsivulle. Tiivistelmässä tulisi ilmetä
mm.  tutkielman aihe tutkimuksen kohde, populaatio, alue ja tarkoitus
käytetyt tutkimusmenetelmät (mikäli tutkimus on luonteeltaan teoreettinen
ja tiettyyn kirjalliseen materiaaliin, on mainittava tärkeimmät lähdeteokset;
mikäli on luonteeltaan empiirinen, on mainittava käytetyt metodit)
keskeiset tutkimustulokset tulosten perusteella tehdyt päätelmät ja
toimenpidesuositukset.
\end{abstract}

\begin{abstracten}
Second abstract in english (in case the document main language is
not english)
\end{abstracten}



% mandatory
\tableofcontents

% if you want a list of figures
% \listoffigures

% if you want a list of tables
\listoftables

% if you want a list of acronyms
\listofacronyms

% change the name if the default doesn't sound right
\renewcommand{\algorithmname}{\listingscaption}

% The thesis starts here.

\begin{comment}
To better organize things, create a new tex file for each chapter
and input it below.

Avoid using the å, ä, ö or <space> characters in referred names and
underscores \_ in file names (may break hyperref).

Good luck! Thanks.
\end{comment}

\chapter{Johdanto} \label{Johdanto}

Viittaaminen lukuun \ref{Johdanto}, toiseen lukuun \ref{Toinen luku},
alilukuun \ref{Alaotsikko}, tätä alempaan lukuun \ref{Alempiotsikko},
alimpaan lukuun \ref{Alinotsikko}, kuvaan \ref{Kuva esimerkki} ja
tauluun \ref{Tauluesimerkki}.

Kuva liitetään seuraavasti. ShareLaTeXin autocomplete rakentaa koko
begin-end blockin yleensä puolestasi.

\begin{figure}
\centering \includegraphics[width=0.5\textwidth]{kuvat/utulogoen}
\caption{Kuvan otsikko}
\label{Kuva esimerkki} 
\end{figure}

Taulukkoja tehdään seuraavasti.

\begin{table}
\begin{centering}
\caption{Taulukon otsikko tulee taulun yläpuolelle}
\begin{tabular}{l|c|r|}
Taulun  & elementit  & erotetaan \tabularnewline
\hline 
toisistaan  & et-merkillä  & \tabularnewline
soluja voi myös  &  & jättää tyhjäksi \tabularnewline
\end{tabular}
\par\end{centering}
\centering{}\label{Tauluesimerkki}
\end{table}

Kirjallisuusviitteet lisätään bib-muodossa bibliografia tiedostoon
ja niihin viitataan niiden ID:llä, joka on bib-muodon ensimmäinen
kenttä \cite{crawley2007write}.

\section{Alaotsikko}

\label{Alaotsikko}

Uskonpuhdistuksen myötä suomi tuli koko jumalanpalveluksen kieleksi.
Raamattu ja liturgiset kirjat oli siksi saatava suomeksi. Maahan tarvittiin
suomea osaavia pappeja; kouluihin piti sen vuoksi saada suomen kielen
opetusta, ja sitä varten tarvittiin oppikirjoja. Nämä asiat olivat
nuoren Mikael Agricolan kannustimena, kun hän aloitti elämäntyönsä
suomen kirjakielen kehittäjänä.

Agricola opiskeli monen muun suomalaisnuorukaisen tavoin Wittenbergissä.
Jo ennen Wittenbergin vuosia hän oli saanut valmiiksi Abckirian ja
Rucouskirian. Abckiria oli tarkoitettu oppikirjaksi. Se sisälsi aakkoset,
tavausharjoituksia ja katekismuksen. Laajassa Rucouskiriassa on rukousten
lisäksi Raamatun tekstejä, muun muassa 41 psalmia. Alussa on monipuolinen
kalendaario, joka sisältää esimerkiksi ruokailu- ja terveydenhoito-ohjeita
ja jopa jonkinlaisen horoskoopin.

Uutta testamenttia Agricola käänsi Wittenbergissä apuneuvoinaan kaksi
latinalaista, kaksi saksalaista ja kaksi ruotsalaista käännöstä. Se
Wsi Testamenti ilmestyi 1548. Kirjan sanasto ja muoto-oppi on siinä
määrin epäyhtenäistä, että on arveltu, että käännöksellä on Agricolan
lisäksi ollut myös muita viimeistelijöitä.

Agricolan osuus 1551 ilmestyneen Psalttarin psalmisuomennoksista on
epäselvä. Suuri osa psalmeista onkin todennäköisesti suomennettu Turun
koulussa Paavali Juustenin johdolla. Juusten itse on kirjoittanut
Psalttarista Suomen piispainkronikassa (suom. Simo Heininen): ``Mutta
ei ole ollenkaan väliä, kenen nimissä se on julkaistu, sillä se on
käännetty, jotta siitä olisi suurta hyötyä Suomen kansalle.'' Pääosa
Psalttarin esipuheista on Agricolan omaa tekstiä. Runomuotoiseen esipuheeseen
sisältyy ansiokas luettelo suomalaisten pakanallisista jumalista.
Agricola suomensi myös osia Mooseksen kirjoista ja profeetoista. Hänen
nimissään on ilmestynyt suomeksi noin 2/5 Raamatusta. Toinen esimerkki
viittaamisesta, jossa myös cite komennon tagi löytyy Bibliografia.bib
tiedostosta \cite{puasuareanu2009survey}.

\subsection{Alempiotsikko}

\label{Alempiotsikko}

Lorem ipsum dolor sit amet, consectetur adipiscing elit. Etiam eget
tellus porttitor, tempus lacus non, pellentesque ligula. Donec sit
amet erat condimentum, feugiat mi accumsan, euismod quam.

Mauris laoreet maximus aliquet. Mauris at gravida elit. Ut nec lobortis
elit. Sed lacinia nisi in ex sollicitudin, ac consequat lacus imperdiet.
Etiam et velit eu lacus maximus faucibus.

\subsubsection{Alinotsikko, joka ei näy sisällysluettelossa}

\label{Alinotsikko}

Lorem ipsum dolor sit amet, consectetur adipiscing elit. Etiam eget
tellus porttitor, tempus lacus non, pellentesque ligula. Donec sit
amet erat condimentum, feugiat mi accumsan, euismod quam.

\paragraph{Otsikko tekstissä, joka ei näy sisällysluettelossa}

Mauris laoreet maximus aliquet. Mauris at gravida elit. Ut nec lobortis
elit. Sed lacinia nisi in ex sollicitudin, ac consequat lacus imperdiet.
Etiam et velit eu lacus maximus faucibus. Vestibulum ante ipsum primis
in faucibus orci luctus et ultrices posuere cubilia Curae; Donec vulputate
tellus ullamcorper odio sodales, non scelerisque neque eleifend. 

\chapter{Raportointi ja PLM järjestelmät} \label{Raportointi ja PLM järjestelmät}

\section{PLM-strategia ja PLM-järjestelmät lyhyesti} \label{PLM-järjestelmät}

Laajempana käsitteenä tuotteen elinkaaren hallinta eli PLM voidaan nähdä yrityksen strategiana hallita tuotetietoja. PLM strategiana koostuu tuotteista, organisaatioista, työmenetelmistä, prosesseista, ihmisistä ja lopulta usein myös tietoteknisestä elinkaaren hallintajärjestelmästä.
Tuotteen elinkaari voidaan jakaa alku-, keski- ja loppuvaiheeseen. Tuotteen elinkaaren pääpiirteet on hyvä ymmärtää, jotta PLM-käsitettä voidaan tarkastella syvällisemmin. Bouhaddoun (2012) konferenssiartikkeli "\textit{PLM Model for Supply Chain Optimization}" määrittelee tuotteen elinkaaren vaiheet ja niiden piirteet seuraavasti: \cite{bouhaddou_plm_2012}

\begin{itemize}
	\item Alkuvaiheessa tuotteen vaatimuksia määritellään ja tuote on luonnosvaiheessa. Luonnosvaiheessa tuotetta voidaan kutsua prototyypiksi (engl. \textit{prototype}) tai mallinnukseksi (engl. \textit{mockup}).
	\item Keskivaiheessa tuote siirtyy tuotantoon ja valmistukseen. Tässä vaiheessa toteutetaan laadunvalvontaa ja kasaamista, ja voidaan puhua jo varsinaisesta tuotteesta. Valmis tuote siirtyy jakemluverkoston kautta itse asiakkaalle. Kun tuote on asiakkaalla, korostuu tuotteen käyttö sekä mahdollinen huolto ja asiakastuki.
	\item Loppuvaiheessa tuotteen elinkaari päättyy. Tuotteen valmistusta ei koeta enää tarpeelliseksi, joten tässä vaiheessa huomio keskittyy tuotannon lopettamiseen ja tuotteen kierrätykseen.
\end{itemize}

Alemanni, ym. (2008) esittää PLM:n suorituskyvyn analysointia käsittelevässä artikkelissaan, että PLM-strategian keskittyessä olennaisesti tuotedatan hallintaan, on PLM-ohjelmisto olennainen osa strategian hyödyntämistä käytännössä. \cite{alemanni_key_2008} Alemanni korostaa, että PLM-järjestelmän kehittäjän tulee kuitenkin toimia yhteistyössä asiakkaiden kanssa, jotta tuotteiden elinkaaren eri vaiheet ja prosessit voidaan implementoida osaksi ohjelmiston toimintoja siten. PLM-käsitteeseen liittyvien määritelmien lisäksi on tärkeää ymmärtää PLM-strategian ja -ohjelmistojen hyötyjä, jotta niiden hyödyntämisen motiivit voidaan ymmärtää. Tarkoituksena on siis vastata siihen, miksi ylipäätään PLM-järjestelmiä käytetään ja kehitetään.

\subsection{PLM-strategian hyödyt ja merkitys} \label{PLM strategian hyödyt}

PLM-strategian hyötyjä on käsitelty laajasti \cite{alemanni_key_2008} \cite{rivest_product_2012}. Strategian hyödyt voidaan jakaa kahteen osa-alueeseen: lyhyen ja pitkän aikavälin hyötyihin. Tietoteknisten PLM-järjestelmien tarkoituksena on taas mahdollistaa PLM-strategian käyttöön ottaminen ja hyödyntäminen käytännössä koko yrityksen tasolla. Näiden järjestelmien pääasiallisena tarkoituksena on koota tietoa yrityksen tuotteiden koko elinkaaren vaiheista keskitettyyn tietojärjestelmään. Tämä mahdollistaa laajojen tuotekantojen johdonmukaisen ja keskitetyn hallinnan yhteistyössä yrityksen eri osastojen ja kumppaneiden välillä.
Konkreettisia PLM:n hyötyjä voidaan havaita Lee, ym. (2008) toteuttamasta tutkimuksesta PLM:n hyödyntämisestä ilmailualalla: IBM-Dassaultin PLM-järjestelmää käytettiin lentokoneiden elinkaaren hallinnassa, jolloin ohjelmiston hyödyntäminen vähensi valmistusaikaa 16:sta kuukaudesta seitsemään kuukauteen. Lisäksi Teamcenter PLM -järjestelmä laski tuotantosykleihin käytettyä aikaa 35:llä prosentilla ja valmistusaikaa 66:lla prosentilla. Keskitetyssä järjestelmässä myös lentokoneiden huoltotarve voitiin ottaa paremmin huomioon jo suunnitteluvaiheessa, mikä suoraviivaisti myös tuotteiden huoltoa niiden elinkaaren aikana. \cite{lee_product_2008}

\subsubsection{Lyhyen aikavälin hyödyt} \label{Lyhyen aikavälin hyödyt}

Lyhyellä aikavälillä PLM-strategia ja PLM-järjestelmän käyttöönotto voi vähentää aikaa jota käytetään työntekijöiden jokapäiväisten työtehtävien suorittamiseen. Strategian avulla yrityksen tuotetiedot ovat keskitetysti saatavilla, eikä ajantasaisia tietoa tarvitse kysellä eri osastojen välillä. Tämä johtaa siihen, että työntekijät voivat käyttää enemmän aikaa tehtäviin, jotka tuottavat yritykselle lisäarvoa arvoa. Lisäksi tuotteiden rakenteiden ymmärtäminen ja visualisointi helpottuu PLM-järjestelmän käyttöönoton myötä. Rakenteen ymmärrystä ja jaettavuutta eri osastojen välillä voidaan parantaa entisestään myös PLM-järjestemän raportoinnilla. \cite{alemanni_key_2008}

\subsubsection{Pitkän aikavälin hyödyt} \label{Pitkän aikavälin hyödyt}

Pidemmällä aikavälillä hyödyt näkyvät konkreettisemmin PLM-strategiaa hyödyntävien yritysten tunnusluvuissa, erityisesti myyntikatteessa. PLM-järjestelmien keskeinen hyöty on prosessien suoraviivaistaminen, mikä johtaa usein tuotteiden saamiseen nopeammin ja useimmin markkinoille. Kun tuotteet pääsevät nopeammin suunnittelusta markkinoille, niiden suunnitteluun ja kehittämiseen käytetyt kustannukset laskevat. \cite{bouhaddou_plm_2012} \cite{alemanni_key_2008}

\section{Raportointi PLM-järjestelmässä} \label{Raportointi PLM-järjestelmässä}

Kuten osioissa \ref{PLM strategian hyödyt} todettiin, yksi PLM-järjestelmän hyödyistä on yksittäisten työntekijöiden työmäärän vähentäminen ja prosessien suoraviivaistaminen. Tietoteknisten järjestelmien etuna on varsinkin automoitu laskenta, joka voi vähentää inhimillisiä virheitä. \cite{niu_organizational_2021}. \cite{rakovic_digital_2022} Tätä automoitua laskentaa voidaan erityisesti hyödyntää raporttien muodostamisessa.
Koska PLM:n tarkoituksena mahdollistaa koko tuotantoketjun yhteistyön asiakkaiden, kehittäjien, toimittajien ja valmistajien välillä tuotteen eri elinkaaren vaiheissa, \cite{bouhaddou_plm_2012} on tärkeää että tuotetieto elinkaaren eri vaiheissa on dokumentoitavissa, analysoitavissa ja helposti jaettavissa. Vaikka PLM-järjestelmien tietomallit ovat usein ohjelmistokohtaisia ja harvemmin standardinomaisia \cite{SritiMohamed-Foued2012PTaS} on tärkeää, että tuotedataa on mahdollista viedä järjestelmän ulkopuolelle tallennettavaksi ja jaettavaksi.

PLM:n kontekstissa raporteilla tarkoitetaan tuotedataa kokoavia ja analysoivia kokonaisuuksia. Raportit voivat olla esimerkiksi PDF- tai Excel-tiedostoja, jotka kokoavat tuotetietoja ja suorittavat laskentaa visualisoimalla dataa esim. kuvaajin tai interaktiivisia "kojelautoja" (engl. \textit{dashboards}), jotka kokoavat useita kuvaajia ja laskettuja arvoja yksittäiseen käyttäjäystävälliseen näkymään. Tässä tutkielmassa keskitytään enemmän raporttitiedostojen tuottamiseen ohjelmallisesti, mutta usein näihin tiedostoihin on myös mahdollista upottaa kojelautamaisia ominaisuuksia, kuten kuvaajia ja tilastoja.
Raporttitiedostojen tuottamista varten monet PLM-järjestelmät tarjoavat "raportointimoottorin" osana PLM-ohjelmistoa, jonka tarkoituksena on kerätä analysoida dataa kokoamalla sitä dokumenttitiedostoformaatteihin. Alan standardina näistä formaateista raporttien kontekstissa lienee PDF-, Excel- ja HTML-pohjaiset raportit, sillä useimmat raportointimoottorit tarjoavat raportteja näissä tiedostomuodoissa ja ne ovat myös tuttuja suurimmalle osalle ohjelmiston käyttäjistä. \footnote{Ajatuksena oli tehdä raportoinnista oma osio ja tämän kappaleen yläotsikko, mutta aihealueena se on hyvin laaja ja väljä. Ehkä on parempi keskittyä raportointiin ja sen merkitykseen PLM-kontekstissa?}

\subsection{Osaluettelo PLM-järjestelmän sydämenä} \label{Osaluettelo PLM-järjestelmän sydämenä}

Yksi PLM-järjestelmän tärkeimmistä toiminnallisuuksista on tuotteen osaluettelon (BOM\nomenclature[BOM]{BOM}{Bill of Materials, osaluettelo, tuoterakenne})(engl. \textit{Bill of Materials}) esittäminen organisoidusti. \cite{david_what_2016} Yksinkertaisuudessaan osaluettelo on lista kaikista osista, joita tarvitaan tuotteet valmistamiseen. Osaluettelossa jokainen yksittäiseen osaan voidaan liittää useita tietokenttiä kuten valmistaja, versio, materiaali ja määrä. Osaluettelo koostuu usein hierarkkisesti osakokoonpanoista, välikokoonpanoista, osakomponenteista ja yksittäisistä osista, eli se kerää dataa siitä, kuinka eri tuotteen komponentit ovat riippuvaisia toisistaan. Osaluetteloa voidaan käyttää viestintään esimerkiksi valmistuskumppanien välillä tai se voidaan rajoittaa yhteen tuotantoyksikköön. \cite{jones_visualizing_2023}

Koska osaluettelot ovat hyvin olennainen osa PLM-järjestelmää, ovat ne tärkeä kohde myös raportoinnille. Osaluetteloista tuotetut raportit voivat analysoida rakennetta pintaa syvemmältä sekä luoda helposti ymmärrettävän yleiskatsauksen massiivisen osaluettelon omaavaan tuotteeseen tarjoamalla samalla esimerkiksi graafeja ja statistiikkaa tuotteesta. Koska raportit voidaan tuottaa erillisinä sähköisinä dokumentteina, voidaan laskentaa jatkaa esimerkiksi Excel-raporttien tapauksessa, tai erityisesti PDF-raportit ovat omiaan arkistoinnille myöhempää käyttöä varten.

\subsection{Raportointi ja Business Intelligence} \label{Raportointi ja Business Intelligence}

Liiketoimintatiedon hyödyntämisellä (BI\nomenclature[BI]{BI}{engl. Business Intelligence, liiketoimintatiedon hyödyntäminen}), (engl. \textit{Business Intelligence}) tarkoitetaan yrityksen kykyä hyödyntää dataa merkityksellisellä tavalla. PLM:n kontekstissa BI korostuu etenkin tuotetiedon hyödyntämisessä. Tätä PLM:n ja BI:n yhteyttä on tutkinut Bosch-Mauchand, ym. (2014) artikkelissaan "\textit{Preliminary Requirements and Architecture Definition for Integration of PLM and Business Intelligence Systems}" \cite{bayro-corrochano_preliminary_2014}. Bosch-Mauchand totetaa, että PLM järjestelmä kulkee käsikädessä BI:n kanssa ja PLM-järjestelmän tuotedatan integraatio ja sen jaettavuus eri järjestelmien välillä on hyvin tärkeää tuotetiedon merkityksellisen hyödyntämisen kannalta.
Bosch-Mauchand erittelee, että jotkin PLM-järjestelmät tarjoavat erillisiä moduuleja raporttien tuottamiseen, mutta harva raportointityökalu tai -moduuli hyödyntää BI:n periaatteita. Bosch-Mauchandin mukaan varsinkin kahden tyyppisillä raporteilla voidaan tuottaa lisäarvoa: \cite{bayro-corrochano_preliminary_2014}
\begin{itemize}
	\item Dokumenttien ja objektien määrällinen analysointi. Esimerkiksi näiden summien tai tyyppien laskenta.
	\item PLM-järjestelmien ominaisuuksien käyttö. Esimerkiksi tuotteen osien uudelleenkäyttö ja tietokantakyselyt.
\end{itemize}
Näiden lisäksi raportointia voidaan hyödyntää IT-hallinnon osa-alueilla, kuten esimerkiksi järjestelmän suorituskyvyn monitoroinnin kannalta.

\section{Raportointimoottorit} \label{Raportointimoottorit}

Raportointimoottori on ohjelmisto tai osa ohjelmistokokonaisuutta, jolla voidaan luoda ohjelmallisesti raportteja. Tämä tarkoittaaa sitä, että raportointimoottori saa vain syötedatan ja mahdollisen pohjan (engl. \textit{template}), mutta kaikki logiikka näiden yhdistämiseen sekä raportin koostamiseen jää raportoinimoottorin vastuulle. \cite{he_design_2010} Raportointimoottorilla ei tarkoiteta ainostaan ohjelmistokehittäjille käytössä olevaa sovellusta tai ohjelmistokirjastoa, sillä raportointimoottorit ovat usein konfiguroitavissa myös loppukäyttäjien toimesta. Täten raportointimoottoria kehittäessä tulee huomioida myös loppukäyttäjän tarpeet. \cite{adhi_performance_2019}

Bambang Prasetya Adhi ym. vertailevat tutkimusartikkelissaan "\textit{Performance comparison of reporting engine birt, jasper report, and crystal report on the process business intelligence}" (2019) \cite{adhi_performance_2019} suosituimpia kaupallisia (SAP Crystal Reports) ja avoimen lähdekkoodin (BIRT ja Jasper Report) raportointimoottoreita. Ahdi ym. käyttävät kokeellisia menetelmiä mitatakseen kolmea raportointimoottorin osa-aluetta:
\begin{itemize}
	\item Soveltuvuus, esimeriksi kuinka hyvin raportointimoottori tukee erilaisen lähdedata käyttöä
	\item Käytettävyys, joka ilmenee oppimisen helppoutena sekä toiminnalisuuden loogisuutena ja tehokkaana käyttönä
	\item Tehokkuus, joka mittaa itse järjestelmän tehokkuutta, esimerkiksi suoritusaikaa
\end{itemize}
Näitä osa-alueita arvioimalla voidaan tehdä perusteltuja päätöksiä raportointimoottorin valinnasta, joten nämä ovat tärkeitä seikkoja ottaa huomioon uuden raportointimoottorin suunnittelussa ja kehityksessä.

\subsection{Raportointimoottorin rakenne ja prosessi} \label{Raportointimoottorin rakenne ja prosessi}

Prosessina raportointimoottori toimii kolmella tasolla: data-, logiikka- ja esitystasolla.\cite{he_design_2010}

\subsubsection{Datataso}

Datatasolla reportointimoottori voi hakea dataa suoraan tietotokannasta tai esimerkiksi API\nomenclature[API]{API}{ohjelmointirajapinta, engl. Application Programming Interface}:n eli ohjelmointirajapinnan (engl. \textit{Application Programming Interface}) kautta. Raportontimoottorin tapauksessa ohjelmointirajapinta voi olla esimerkiksi hakurajapinta, jonka taustalla toimivan hakumoottorin avulla reportontimoottori voi hakea tarvitsemaansa dataa täsmällisemmin. Varsinkin PLM-järjestemän kontekstissa hakumoottori on hyvin keskeinen osa PLM-järjestelmää. Datataso määrittelee siis mistä ja miten raportointimoottori hankkii lähtödataa sekä millaiset lähtötiedot sillä on koostaa raportti. Näihin lähtötietoihin lukeutuu esimerkiksi mahdolliset raporttipohjat ja muut käyttäjän määrittämät asetukset.

\subsubsection{Logiikkataso}

Logiikkatasolla raportointimoottori jäsentelee lähtödataa ja suorittaa laskentaa. Lähtödatan formaatti on usein historiallisesti ollut XML\nomenclature[XML]{XML}{merkintäkielien standardi ja tiedostomuoto, engl. Extensible Markup Language} (engl. \textit{Extensible Markup Language}) sen ollessa yksi internetin yleisimmin käytetyistä dataformaateista, mutta JavaScriptin yleistyttyä myös JSON\nomenclature[JSON]{JSON}{engl. JavaScript Object Notation, yksinkertainen tiedstomuoto tiedon välitykseen ja tallennukseen} (engl. \textit{JavaScript Object Notation}) noussut suosituksi tiedostomuodoksi. JSONin etuna on sen suora integraatio JavaScriptin yhteyteen sekä sen nopeus verrattuna vanhempaan XML-standardiin. \cite{nurseitov_comparison_nodate} Logiikkatason toteuttaman laskennan avulla lähtödatasta voidaan luoda esimerkiksi graafeja ja sekä taulukoida dataa, mikä mahdollistaa lasketut sarakkeet ja summat. Tämä voidaan nähdä raportointimoottorin ytimenä, sillä se tuottaa merkityksellistä dataa usein vaikeaselkoisesta lähdedatasta. Lisäksi logiikkatason tulee jäsentää data siten, että se on mahdollista kirjoittaa tiedostoon esitystasolla.

\subsubsection{Esitystaso}

Esitystasolla data kirjoitetaan tiedostoon, jolloin se voidaan tallentaa käyttäjäystävällisessä tiedostoformaatissa. Suosituimpia tiedostoformaatteja ovat HTML-, Excel- ja PDF-tiedostoformaatit. Esitystasolla raportin merkitys ilmenee käyttäjälle: lähtödata on esitetty helposti omaksuttavassa ja ymmärrettävässä muodossa, sekä raportti on luettavissa ja jaettavissa helposti yksittäisenä tiedostona.

\subsection{PLM-järjestelmä raportointimoottorin toimintaympäristönä}

Kuten kappaleessa \ref{Osaluettelo PLM-järjestelmän sydämenä} todettiin, osaluettelot ovat keskeisesen osa PLM-järjestelmään tallennetun datan esittämistä. Koska tuoteobjektit koostuvat osaluetteloista, myös PLM-järjestelmässä tuotteista koostettavat raportit perustuvat osaluetteloista kerättyyn lähtödataan. PLM-järjestelmän raportointimoottorit ovat siten erikoistuneita jäsentelemään ja kokoamaan hierakkista dataa. PLM-järjestelmän tarjoamille raporteille on olennaista tuotteeseen ja sen kehitykseen liittyvät seikat, kuten esimerkiksi tuotteen osien toimittajien jakauma tai tuotteen muokkaushistoria. Lisäksi osaluettelon perusteella voidaan laskea yksittäisten osien summia rakenteessa tai esimerkiksi luoda raportteja tuotteen osista, jotka täyttävät tietyt kriteetrit.

Mahdollisen raportointimoottorin ohella PLM-järjestelmät sisältävät usein hakumoottorin, jonka avulla voidaan etsiä tehokkaan ja tarkan haun tietokannasta annettujen kriteerien mukaisesti. Raporttien muodostamisessa ulkoisen hakumoottorin hyödyntäminen vähentää itse raportontimoottorin kuormaa, jolloin raportontimoottorin toiminallisuuden kehittämisessä voidaan keskittyä enemmän laskentaan ja lisäarvon tuottamiseen. Täten PLM-järjestelmän tapauksessa lähtödatan hakeminen voi tapahtua haku-API:n välityksellä, jolloin raportteja voidaan muodostaa tuotterakenteiden lisäksi esimerkiksi jonkin tietyn hakulausekkeen perusteella.

PLM-järjestelmien käyttäjät ovat tyypillisesti suhteellisen suuren mittakaavan teollisuusyrityksiä. Useissa tapauksesissa myös tuoterakenteet ovat valtavia, joten raportoinimoottorin tulee olla tarpeeksi tehokas ja optimoitu, jotta myös suurista tietorakenteista on mahdollista koostaa raportteja siedettävässä suoritusajassa. Raportointimoottorin logiikkatason lisäksi PLM-järjestelmien käyttäjillä on myös tarpeita muokata raporttien ulkoasua raportointimoottorin esitystasolla. Esimerkiksi yrityksen logojen ja visuaalisen ilmeen muokkaaminen raportteihin on olennainen osa raportontimoottorin toiminnallisuutta.

\chapter{Tapaus: Sovelia PLM:n raportointityökalu} \label{Tapaus: Sovelia PLM:n raportointityökalu}

\section{Sovelia PLM} \label{Sovelia PLM}

Kehitettävän raportointityökalun toimintaympäristönä toimii kaupallinen PLM-järjestelmä, Sovelia PLM. Sovelia PLM:n kehitys on alkanut yli 30 vuotta sitten ja sen kehitys jatkuu aktiivisesti edelleen. Kuten muutkin PLM-järjestelmät, Sovelia PLM pyrkii ratkaisemaan valmistusalan yritysten haasteita liittyen tuotteen datan hallintaan sen koko elinkaaren ajan. \cite{soveliaAboutSovelia} Sovelia PLM:n erityispiirteenä on sisältämät valmiiksi konfiguroidut \textit{"templatet"} eli valmiit mallit objekteille ja prosesseille, jotka ovat muokattavissa asiakkaan tarpeiden mukaan. Lisäksi malleihin kuuluu muita valmiiksi konfiguroituja työkaluja että alalla hyväksi todettuja prosesseja. \cite{soveliaSoveliaGetting}

Sovelia PLM koostuu objekteista ja objektilinkeistä. Objekteja voidaan määritellä niiden attribuuttien avulla. Objektit voivat olla osia, piirustuksia, dokumenttilinkkejä tai linkkejä muihin toimintoihin. \cite{soveliaSoveliaGetting} Objektilinkkien avulla voidaan muodostaa osaluetteloita, jotka ovat tärkeitä raportoinnin kohteita. Näiden lisäksi toinen merkittävä konsepti Soveliassa on käyttäjä ja käyttäjäryhmät. Yksinkertaisuudessaan järjestelmällä voi olla luonnollisesti useita käyttäjiä ja käyttäjät voivat kutsua eri käyttäjäryhmiin. Käyttäjäryhmiä voi olla useita, mutta tärkein niistä on ymmärtää \textit{"admin"}-ryhmä (engl. \textit{administrator}), johon kuuluu pääkäyttäjät eli järjestelmänvalvojat. Järjestelmänvalvojalla on luonnollisesti oikeuden muuttaa järjestelmän asetuksia. Pääkäyttäjän konsepti on tärkeä ymmärtää, sillä osa kehitettävän raportointityökalun ominaisuuksista on saatavilla vain pääkäyttäjälle. Käyttäjäryhmien lisäksi jokaisella käyttäjällä on lisenssi, joka määrittelee osaltaan käyttäjän oikeuksia rajoittamalla esimerkiksi objektien ja objektilinkkien muokkausoikeuksia. \cite{soveliaUserManagement}

\subsubsection{Raportointi Sovelia PLM-järjestelmässä} \label{Raportointi Sovelia}

Sovelia PLM ja muut PLM-järjestelmät poikkeavat valmiiden raportointityökaluratkaisujen tyypillisistä toimintaympäristöistä, sillä kuten luvussa \ref{PLM-järjestelmä toimintaympäristönä} huomattiin, PLM-data on usein kompleksista ja monimuotoista. Siemensin PLM-järjestelmän verkkosivuilla julkaistussa artikkelissa \textit{"The Challenge of Getting High Quality Reports out of PLM"} \cite{german_challenge_2016} (2016) esitellään haasteita, joita liittyy korkealaatuisten raporttien tuottamiseen PLM-järjestelmistä. Artikkeli nostaa esille osaluetteloiden (BOM) merkityksen PLM-järjestelmän raportoinnissa: on tärkeää että raporttien saama ja tuottama data on luotettavaa ja laadukasta, jotta päätöksenteko raporttien pohjalta olisi mahdollista. Osaluetteloihin voi tulla jopa satoja muutoksia päivittäin, joten raporttien tapauksessa on tärkeää, että niiden käyttäjät tietävät työskentelevänsä oikean datan kanssa. Täten raporttien ajantasaisuus ja selkeät aikaleimat ovat hyvin tärkeitä PLM-järjestelmän ja myös Sovelian tuottamille raporteille. Sovelia PLM:n raportointimoottorin etuna on hakurajapinnan käyttö, joka palauttaa aina hakupyynnön mukaisesti ajankohtaista dataa, mikä vähentää mahdollisten virheellisten tietojen määrää.

PLM-datan kompleksisuus tekee raportointityökalun kehittämisestä haasteellista, sillä myös Sovelia PLM vaatii erityisesti tarkoitusta varten kehitetyn raportointityökaluratkaisun. Tarkoitusta varten kehitetty ratkaisu on ohjelmistokehittäjälle työläs rakentaa alusta asti itse, mutta sen etuna on sen täysi muovautuvuus tarvetta varten. On kuitenkin todettava, että myös valmiin raportointimoottorin ja -työkalun implementoiminen Sovelia PLM:n toimintaympäristöön olisi todennäköistä työlästä juuri PLM-datan kompleksisuuden vuoksi.

\subsubsection{Sovelia PLM:n vanha raportointityökalu}

Sovelia PLM -järjestelmässä on tuotannossa ja asiakkailla käytössä vanha raportointiyökalu, joka tarjoaa raportteja PDF, Excel ja ZIP-tiedostoformaateissa. Näistä poikkeavimpia ovat ZIP-raportit eli ZIP-arkistot, \nomenclature[ZIP]{ZIP}{tiedon pakkaukseen käytetty tiedostomuoto, joka voi sisältää useita tiedostoja ja kansioita pakattuna} jotka sisältävät useita PDF- tai Excel-raportteja pakattuna yhteen tiedostoon. Täten ZIP-raporttiin voidaan sisällyttää esimerkiksi raportti tuoterakenteen jokaiselle yksittäiselle objektille, jolloin ZIP-paketin kansiorakenne vastaa itsessään tuotteen rakennetta. Raporttidokumenttein data kerätään \textit{"rakenneagentin"} avulla, joka on nimensä mukaisesti osaa kulkea objekti-linkki -suhteita pitkin ja täten kerätä tarvittavan lähdedatan raportin koostamista varten. Itse raporttidokumentteihin kirjoittaminen tapahtuu Java-ohjelmointikielen Apache FOP-ohjelmakirjaston \cite{noauthor_apachetm_nodate} avulla XSL-\nomenclature[XSL]{XSL}{XML-kieliperhe, joka mahdollistaa XML-pohjaisten tiedostojen ulkoasun ja rakennemuutoksen määrittelyn} (engl. \textit{Extensible Stylesheet Language}, kieliperhe, joka mahdollistaa XML-pohjaisten tiedostojen ulkoasun ja rakennemuutoksen määrittelyn) ja XML-tietorakenteisiin tallennetun datan antaman ohjeistuksen avulla.

Vaikka tämä ratkaisu on toimiva ja edelleen kelvollinen tekniikka dokumenttien muotoiluun, ja siten raporttien generointiin ohjelmallisesti ennalta määritellyn datan mukaisesti, se kohtaa ongelmia datamäärien kasvaessa. Näihin ongelmiin kuuluu esimerkiksi liiallinen muistinkäyttö palvelimella, prosessin hitaus varsinkin suurempien tuoterakenteiden tapauksissa ja raporttien konfiguroinnin haasteellisuus sen ollessa täysin mahdotonta PLM-järjestelmän loppukäyttäjille. Lisäksi toiminnon ylläpito alati muuttuvassa toimintaympäristössä muuttuu todennäköisesti haasteellisemmaksi tulevaisuudessa, kun uusi teknologia syrjäyttää vanhaa yhä useammilla osa-alueilla, samalla kun vanhalle pohjalle rakennettu raportointimoottori ei pysy samalla tavalla muutoksen tahdissa.

\subsubsection{Uuden raportointityökalun haasteet} \label{Uuden haasteet}

Koska vanhan raportointityökalun modernisointi ja päivittäminen nykypäivään vaatisi perustavanlaatuisia muutoksia koko järjestelmään, päädyttiin valitsemaan kehityskohteeksi kokonaan uuden raportointityökalun kehittämisen. Uuden työkalun tarkoituksena on integroitua paremmin nykyiseen Web-pohjaiseen käyttöliittymään ja tarjota käyttäjäystävällisemmän kokemuksen raporttien luomiseen. Raporttien luominen tuoterakenteista tapahtuisi käyttäjän näkökulmasta paremmin integroituna Web-käyttöliittymään tarjoten mahdollisuuden seurata meneillään olevien raporttien edistymistä sekä perua raporttien koostamisprosesseja käyttöliittymästä käsin. Lisäksi pääkäyttäjällä tulisi olla mahdollisuus muokata ja luoda uusia raporttityyppejä muokkaamalla esimerkiksi raporteissa esitettäviä sarakkeita ja tuoterakenteiden suodattamiseen liittyviä sääntöjä. Myös raporttien ulkonäön tulee olla muokattavissa pääkäyttäjän toimesta, mikä tarkoittaa niiden visuaalisen ilmeen muokkausta lisäämällä asiakasyrityksen logoja, muokkaamalla väriteemaa ja asettelua asiakkaan tarpeiden mukaisesti.

Uuden raportointityökalun tulee ensimmäisessä vaiheessa sisältää samat toiminnallisuudet kuin vanha raportointityökalu, mutta tarkoituksena on tuottaa raportteja huomattavasti nopeammin aiheuttamalla samanaikaisesti vähemmän kuormaa palvelimelle. Raporttien sisällön ja yleisen rakenteen tulee vastata vanhempia raportteja identtisesti, mutta raporttien visuaalinen ilme tulee päivittää nykypäivän standardien mukaiseksi. Lisäksi vaaditaan syvällisiä teknologiatutkimusta ja -analyysiä olemassa olevien raportointityökalujen ja -moottorien ratkaisuista, Sovelia PLM:n ohjelmakannassa jo olemassa olevista komponenteista ja muista mahdollisesti hyödyllisistä teknologioista.

\section{Nykyiset erilliset raportointityökalut} \label{Nykyiset erilliset}

Ennen käytettävien teknologioiden ja tekniikoiden valintaa on hyvä perehtyä jo olemassa oleviin raportointityökaluihin ja -moottoreihin, jotta voidaan muodostaa kokonaiskuva yleisesti käytössä olevista ratkaisuista. Hyvän kokonaiskuvan avulla voidaan tehdä valistuneita päätöksiä siitä, miten raportointityökalua kannattaa lähteä kehittämään sekä ymmärtää käytössä olevien tekniikoiden hyötyjä ja rajoituksia PLM-järjestelmän kontekstissa. Analysoimme kuutta suosittua itsenäistä (engl. \textit{standalone}) raportointityökalua, joista neljä ovat kaupallisia ja kaksi ovat avoimen lähdekoodin ohjelmistoja. Juuri nämä kuusi raportointityökalua valittiin tarkasteltavaksi, sillä ne ovat kaikki paikallisesti asennettavia ohjelmistoja (engl. \textit{On-premise software}) kuten myös Sovelia PLM. Täten tässä tutkimuksessa ei keskitytä esimerkiksi ulkoisten pilvipalveluiden raportointityökalu verkko-ohjelmointirajapintoihin (engl. \textit{Web API}). Yksi esimerkki tällaisesta verkkopohjaisesta API:sta on \textit{Bold Reports} \cite{noauthor_bi_2022}, joka tarjoaa monipuolisen verkkorajapinnan kautta toimivan raportointityökalun. Toinen kriteeri tarkasteluun valinnalle oli käytettävä teknologia, sillä Sovelia PLM:n nykyinen arkkitehtuuri nojaa vahvasti Java- ja JavaScript-ohjelmointikieliin, joten täysin uuden ohjelmointikielen lisäämistä järjestelmäarkkitehtuurin pinoon ylläpidettäväksi ei pidetty mielekkäänä vaihtoehtona. Lisäksi esimerkkejä sekä Java- että JavaScript-pohjaisia raportointityökaluista löytyi runsaasti, joten uuden teknologian valitseminen ei vaikuttanut välttämättömyydeltä. Kolmas tärkeä kriteeri oli ulostulodataformaattien monipuolisuus: koska Sovelia PLM -järjestelmään kehitettävän raportointityökalun tulee tuottaa raportteja eri tiedostoformaateissa, tulee myös tarkasteltavien työkalujenkin.

Puhtaasti kaupallisia ohjelmistoja ovat JavaScript-pohjainen \textit{ActiveReportsJS}, SAP:in, maailman suurimman yritysohjelmistoja tarjoavan yrityksen \cite{noauthor_sap_nodate-3}, kaksi raportointityökalua \textit{SAP Business Objects} ja \textit{SAP Crystal Reports} ja Oraclen, talouslehti Forbesin mukaan vuoden 2023 80. suurimman julkisen yrityksen \cite{tucker_global_nodate}, raportointityökalu \textit{Orancle BI Publisher}. Vahva esimerkki taas avoimen lähdekoodin vaihtoehdosta on \textit{JasperReports}, joka kutsuu itseään suosituimmaksi avoimen lähdekoodin Java-raportointikirjastoksi. \cite{noauthor_jasperreports_nodate}. \textit{JasperReports} sisältää myös visuaalisen suunnittelutyökalun, \textit{iReport Designerin}, joten \textit{JasperReports} on kokonainen raportointityökaluohjelmisto. \textit{JasperReports} sisältää myös vaihtoehtoisen maksullisen lisenssin yrityskäyttöön tarjoten yritysasiakkaille parempaa tukea ja ratkaisua suuremman skaalan raportointiin. \cite{noauthor_jasper_nodate} Toinen avoimen lähdekoodin raportointityökalu on myös Java-pohjainen \textit{BIRT} (lyhenne engl. \textit{"Business Intelligence Reporting Tool"}), joka tarjoaa monipuolisia raportointimahdollisuuksia täysin avoimen lähdekoodin ja Eurooppalaisen \textit{Eclipse Foundationin} ylläpitämä ohjelmisto. \cite{eclipsebirtBIRT} \textit{BIRT} on keskittynyt raportoinnin lisäksi myös Business Intelligencen hyödyntämiseen asiakas- ja web-applikaatioissa, joten se sisältää monia erilaisia visualisointimahdollisuuksia ja raportointimoottorin lisäksi myös graafisen raporttikonfigurointityökalun.

\subsubsection{Nykyisten raportointityökalujen yleiskuva} \label{Nykyisten yleiskuva}

Raportointityökalua kehittääkseen osaksi PLM-järjestelmää varten tulee lisenssien lisäksi ottaa huomioon myös muita ominaisuuksia. Liitteenä olevassa taulukossa \ref{tab:Yleista} luokitellaan raportointiyökalujen yleisiä ominaisuuksia. Taulukoiden data on kerätty taulukossa mainituista lähteistä, kuten ohjelmiston verkkosivuilta, ohjelmakirjaston dokumentaatiosta tai myyntiesitteistä. Jo mainitut lisenssit ovat eritelty kaupallisiin ja avoimen lähdekoodin ohjelmistoihin. Avoimen lähdekoodin ohjelmistot tarjoavat asiantuntijoille lähdekoodin perehtymisen avulla syvän katsauksen raportointityökalun sisäiseen toiminnallisuuteen, mutta tässä tutkimuksessa keskitymme enemmän pintapuolisempiin ominaisuuksiin ja teknologiavalintoihin tarkemman lähdekoodianalyysin sijasta. Yksi tällaisista ominaisuuksista on lähdedatan formaatti ja mahdolliset datalähteet. Kaikki tutkitut raportointityökalut tukivat jotakin rakenteellisen datan tallentamiseen suunniteltua tiedostoformaattia eli tässä tapauksessa XML:ää tai JSON:ia. Ainoastaan \textit(ActiveReportsJS) ei tukenut XML:ää, mutta se tukee kuitenkin JavaScript-pohjaisena luonnollisesti JSON:ia. Rakenteellisen datan tukeminen tärkeää varsinkin PLM:n datan tapauksessa, jossa objektien ja siten tuotteiden väliset suhteet ovat avainasemassa. Näiden tiedostoformaattien lisäksi noin puolet tarkasteltavista raportointityökaluista tuki myös SQL-kyselyitä tietolähteenä \nomenclature[SQL]{SQL}{engl. \textit{Structured Query Language}, standardoitu kyselykieli, jolla relaatiotietokantaan voi tehdä erilaisia hakuja, muutoksia ja lisäyksiä} (engl. \textit{Structured Query Language}). Näiden lisäksi myös tavallisten tekstitiedostojen käyttö tietolähteenä on tyypillistä XML- ja JSON-tiedostojen lisäksi, kuten myös Excel-taulukkolaskentaohjelman tuottamat XLSX-tiedostot. Mielenkiintoista on, että myös Excelin tuottamat XLSX-tiedostot koostuvat itseasiasiassa vain ZIP-pakatuista XML-tiedostoista. \cite{miner_under_2022}

Raportointityökalujen ohjelmointikieliä tarkastellessa suosituimmaksi ilmeni Java, sillä kaikki paitsi \textit(ActiveReportsJS) olivat Java-pohjaisia. Yksi syy tähän voi olla Javan suuri ja vakaa suosio 1990-lopussa ja 2000-luvun alussa ohjelmointikielten suosiota mittaavan \textit{TIOBE indexin} mukaan\cite{noauthor_tiobe_nodate}, jolloin suurinta osaa näistä raportointityökaluista alettu kehittää. Esimerkiksi \textit{JasperReportsia} on alettu kehittää vuonna 2001 \cite{noauthor_origin_2006}, \textit{BIRTiä} vuonna 2004 \cite{noauthor_eclipse_2004} ja \textit{BusinessObjectsia} jo vuonna 1995 ennen kuin se kuului SAP-tuoteperheeseen \cite{alma9922048891805971}. Kuten {TIOBE indexistä} voidaan havaita, on Java edelleen suhteellisen suosittu kieli käyttää vuonna 2023, mutta sen suosio on kuitenkin jo hieman laskussa. Raporttipohjien konfigurointi tapahtuu raportointityökaluissa pääsääntöisesti jonkin käyttöliittymän avulla. Käyttöliittymän lisäksi \textit{ActiveReportsJS} tarjoaa mahdollisuuden konfiguroida raporttipohjia JSON-datan avulla, \textit{BIRT} ja  \textit{JasperReports} XML-datan avulla ja \textit{Oracle BI Publisher} RTF (engl. \textit{Rich Text Format}, eroaa tavallisesta tekstistä tallentamalla tiedostoon myös fontit, fonttikoot ja yleisen muotoilun, mm. kursivoinnin) ja XSL-pohjaisesti. Tästä voidaan havaita, että visuaalinen käyttöliittymä raporttipohjien luomiselle on suosituin tapa toimia, mutta tiedostopohjainen raporttipohjien muokkaus on käypä vaihtoehto.

\subsubsection{Nykyisten raportointityökalujen ulostuloformaatit}

Raportointityökalun tärkein tehtävä on tuottaa raportteja, jolloin raporttien tulee olla myös helposti avattavissa ja esitettävissä eri laitteilla. Liitteenä olevassa taulukossa \ref{tab:Formaatit} tarkastellaan samojen kuuden raportointityökalun tarjoamia ulostuloformaatteja, eli sitä millaisissa tiedostomuodoissa työkalut tarjoavat raportteja. Ulostuloformaattien valinta on erityisen tärkeää raportointityökalulle, sillä sen käyttäjät suosivat heille jo entuudestaan tuttuja tiedostomuotoja, sekä ohjelmistoteknisestä näkökulmasta, sillä raporttien tuottaminen erilaisissa tiedostomuodoissa vaatii usein tiedostomuotokohtaisia työkaluja.

Tiedostomuodot voidaan jaotella avoimiin ja suljettuihin tiedostomuotoihin. Avoimessa tiedostomuodossa olevan tiedoston avaaminen on mahdollista ilman kaupallista ohjelmaa, jolloin se on kenen tahansa avattavissa ja sillä on avoimesti saatavilla oleva spesifikaatio. \cite{fileformats} Suljettu tiedostomuoto on luonnollisesti avoimen tiedostomuodon vastakohta, jolloin spesifikaatio ei ole avoimesti saatavilla. Suljettu tiedostomuoto voi olla hyödyllinen joissain tapauksissa, esimerkiksi silloin jos tiedostoon tulee tallentaa tai sen avulla tulee esittää dataa siten että muut esimerkiksi avoimet tiedostomuodot eivät siihen pysty. Avoimien tiedostomuotojen suosiminen on kuitenkin kannattavaa, sillä se vähentää riippuvuutta kolmannen osapuolen ohjelmistoista tai tiedostomuotojen lisenssinhaltijoista. Pahimmassa tapauksessa tietoja voidaan lukea vain tietyillä ohjelmistopaketeilla, jotka voivat olla kohtuuttoman kalliita tai jotka voivat vanhentua. \cite{fileformats}

Ulostuloformaateista erityisen suosittu on Adoben kehittämä PDF (engl. \textit{Portable Document Format}), jota käytetään erityisesti sähköiseen julkaisemiseen. PDF on ollut ISO-standardi vuodesta 2007 alkaen ja ISO (engl. \textit{International Organization for Standardization}, kansainvälinen standardisoimisjärjestö) kuvailee sitä maailman suosituimmaksi dokumenttitiedostoformaatiksi \cite{naden_standard_2021}. Vaikka PDF on alun perin Adoben kehittämä kaupallinen ja suljettu tiedostomuoto, on sen spesifikaatio ISO-standardoinnin jälkeen avoimesti saatavilla. Suosionsa, standardoinnin ja avoimen spesifikaation vuoksi PDF onkin luonnollinen valinta raportointityökalun ulostuloformaatille ja siksi myös kaikki kuusi tarkastelevaa raportointityökalua tarjoavat raportteja PDF-muodossa.

Raporttidokumenteille selkeästi suosittuja alustoja ovat erityisesti yritysmaailmassa suositun Microsoft Officen tuottamat tiedostomuodot, erityisesti taulukkolaskentaohjelma Excelin \textit{.xlsx} ja \textit{.xls} ja tekstinkäsittelyohjelma Wordin tuottamat \textit{.docx} ja \textit{.doc} tiedostot. Microsoft Officen tiedostomuotojen laaja tuki raportointityökalujen saralle ei ole varsinaisesti yllättävää ottaen huomioon, että tilastopalvelu \textit{Statistan} mukaan Microsoftilla on noin 45 prosentin markkinaosuus toimisto-ohjelmistojen alalla Googlella sen ollessa noin 50. \cite{office} On kuitenkin otettava huomioon, että myös Google tarjoaa omissa toimisto-ohjelmistoissaan mahdollisuuden avata ja muokata myös Microsoft Officen tuottamissa tiedostomuodoissa olevia tiedostoja. Sekä Excelin että Wordin tuottamat tiedostomuodot \textit{.xlsx} ja \textit{.docx} ovat suljettuja tiedostomuotoja, mutta ne ovat kuitenkin Googlen tarjoamien työkalujen lisäksi avattavissa esimerkiksi avoimen lähdekoodin ohjelmistolla LibreOfficella \cite{libreoffice}. Täten vaikka Word ja Excel-tiedostot ovat suljetussa tiedostomuodossa, on ne mahdollista avata ilman kaupallisia ohjelmistoja. Microsoft Officen tiedostomuodoista erityisesti taulukkomuotoisten raporttien tapauksessa Excelin \textit{.xlsx}-tiedostot ovat hyödyllisiä, sillä ne tarjoavat raportointityökalun käyttäjälle mahdollisuuden tutkia raportin sisältöä vielä pintaa syvemmältä suorittamalla itse laskentaa ja kokoamalla kuvaajia ohjelman avulla. Raportteja Excel- ja Word-tiedostoformaateissa tarjosivat kaikki tarkasteltavat raportointityökalut paitsi ActiveReportsJS, eli nämä formaatit ovat PDF-raporttien ohella erityisen suosittuja.

Kolmas suosittu tiedostomuoto on HTML (engl. \textit{Hypertext Markup Language}, suom. \textit{hypertekstin merkintäkieli}) \nomenclature[HTML]{HTML}{engl. \textit{HyperText Markup Language}, suom. \textit{hypertekstin merkintäkieli}, merkintäkieli, jolla internetsivut ovat kirjoitettu}, joka on erityisesti verkkosivujen kirjoittamisessa käytetty merkintäkieli. HTML ei ole ohjelmointikieli vaan merkintäkieli, joka määrittelee miten sisältö tulisi muotoilla. Myös HTML on avoin ja ISO-standardoitu tiedostomuoto.\cite{iso_html} HTML:ää käytetään usein yhteydessä CSS:n kanssa (engl. \textit{Cascacing Style Sheets}, suom. \textit{porrastetut tyyliarkit}) \nomenclature[CSS]{CSS}{engl. \textit{Cascacing Style Sheets}, suom. \textit{porrastetut tyyliarkit}, tyylisivu, jolla voidaan määritellä tyyliohjeita miten esimerkiksi HTML dokumentti voidaan esittää}. CSS on eräänlainen tyylisivu, jolla voidaan määritellä tyyliohjeita siitä, miten esimerkiksi HTML dokumentti voidaan esittää. Kuten HTML myöskään CSS ei ole ohjelmointikieli, mutta se tarjoaa valtavasti mahdollisuuksia muotoilla HTML-dokumentteja. HTML:n ja CSS:n ohella verkkosivut hyödyntävät tyypillisesti myös JavaScriptia. HTML-raportit tarkoitettu esitettäväksi verkkoselaimessa, mikä tekee niistä hyvin alustariippumattomia. Mikäli HTML-raportit sisältävät CSS-tyyliarkin ja JavaScriptia, näitä hyödyntämällä HTML raportteihin voidaan lisätä animaatioita ja muita toiminnallisuuksia, ne voivat olla huomattavasti interaktiivisempia käyttäjälle verrattuna staattisempiin PDF-, Excel ja Word-raporttidokumentteihin. HTML-pohjaisia raportteja tarjosivat poikkeuksetta kaikki tarkasteltavat kuusi raportointityökalua.

Viimeisenä ulostuloformaattina vain \textit{SAP BusinessObjects} ja \textit{SAP Crystal Reports} mahdollistavat raporttien tuottamisen näiden raportointityökalujen omassa suljetussa tiedostomuodossa nimeltä \textit{RPT}, joka lyhennelmä englannin kielen sanasta \textit{report}. Koska \textit{RPT} tiedostoformaatti on suljettu ja se ei ole yleisesti laajalti käytössä oleva tiedostoformaatti, \textit{SAP BusinessObjectsin} ja \textit{SAP Crystal Reportsin} ulkopuolella, sen avaaminen on käytännössä mahdollista vain tätä varta vasten kehitetyllä ohjelmalla toisin kuin esimerkiksi Word- ja Excel-tiedostojen tapauksessa.

\section{Uuden raportointityökalun kehitys} \label{Uuden kehitys}

Ensimmäinen haaste raportointityökalun kehittämiselle on käytettävien tekniikoiden ja teknologioiden valinta. Koska tarkoituksena on kehittää täysin uusi toiminnallisuus alusta alkaen, on tärkeää ymmärtää nykyiset Sovelia PLM:ssä käytössä olevat komponentit palvelinarkkitehtuurien näkökulmasta. Itse raportointimoottorin täytyy toimia jollakin palvelimella ja koostaa valmiit raportit palvelinpäässä ennen niiden lataamista käyttäjien saataville. Tällä menettelyllä voidaan vähentää asiakaspään kuormaa siirtämällä raskaampi laskenta palvelimen vastuulle. Palvelinpäässä kuormaa voidaan tasata asettamalla tehtävät jonoon odottamaan suoritusvuoroaan.

\subsubsection{Uuden raportointimoottorin palvelinympäristö}

Tarkastelemalla olemassa olevia raportointityökaluja, Java osoittautui suosituimmaksi ohjelmointikieleksi raportointityökalun ohjelmointirajapinnalle. Luvussa \ref{Nykyisten yleiskuva} kuitenkin todettiin, että Java on suosioltaan laskussa \textit{TIOBE Indexin} mukaan vuonna 2023. Sovelia PLM hyödyntää ohjelmistopinossaan (engl. \textit{software stack}) Java-palvelinta, mutta se on elinkaareltaan enemmänkin ylläpitovaiheessa, eikä sen alaisuuteen ole suunnitelmissa kehittää kokonaan uusia toiminnallisuuksia. Uuden raportointityökalun tarkoituksena on olla mahdollisimman tulevaisuudenkestävä (engl. \textit{future-proof}), joten on olennaista tarkastella myös vaihtoehtoisia palvelinympäristöjä raportointityökalulle.

Raportointimoottorin palvelinympäristöksi valiutui lopulta Sovelia PLM:ssä jo käytössä oleva Node.js. Node.js on avoimen lähdekoodin alustariippumaton ajoympäristö JavaScript-koodin suorittamiseen palvelimella ja se on Stack Overflown teettämän kyselyn mukaan yleisin käytetty Web-teknologioita ammattikehittäjien keskuudessa React.js:n ohella n. 42 prosentin käyttöasteella. \cite{stackoverflowStackOverflow} Kyselyn mukaan Node.js on kuitenkin vuoden 2023 johtava teknologia palvelinpuolella, sillä React.js on tarkoitettu enemmänkin käyttöliittymien kehittämiseen asiakaspäässä. Node.js tarjoaa oletuksena paketinhallintajärjestelmän nimeltään NPM (engl. \textit{Node Package Manager, NPM}, tyyliteltynä npm), joka koostuu komentorivityökalusta ja JavaScript -ohjelmapaketteja ylläpitävästä tietokannasta. NPM:n avulla erillisten ohjelmakirjastojen sisällyttäminen ja käyttöönotto on yksinkertaista ja nopeaa.\cite{npmjsAbout} Node.js:n suosion vuoksi käytettävissä on laaja kattaus erilaisia hyödyllisiä ohjelmakirjastoja. NPM helpottaa erityisesti tarvittavien ohjelmakirjastojen käyttöönottoa ja asennusta, joita raportointimoottori tarvitsee erityisesti raporttidokumenttien kokoamiseen ohjelmallisesti. Sovelia PLM:n Node.js-palvelin tarjoaa erilaisia palveluita ja rajapintoja Sovelia PLM:n Web-käyttöliittymälle. Näihin palveluihin kuuluu esimerkiksi 3D-näkymien ja -tiedostojen tarjoaminen Web-käyttöliittymälle, joten myös raportointimoottorin sisällyttäminen samalle palvelimelle oli luonteva ratkaisu, sillä myös raporttien konfigurointityökalu tullaan kehittämään Sovelia PLM:n web-käyttöliittymään.

\subsubsection{Uuden raportointityökalun ominaisuuksia}

Datalähteenä Sovelia PLM:n uusi raportointimoottori käyttää Node.js pohjaisena JSONia, joka oli käytössä myös muissa tarkastelluissa raportointityökaluissa. Luvussa \ref{Nykyisten yleiskuva} todettiin JSONin olevan hyvä valinta hierarkkisen PLM-datan tallentamiseen sekä luvun \ref{Raportointimoottorin rakenne ja prosessi} alaluvussa koskien logiikkatasoa huomattiin JSONin olevan nopeampi verrattuna vanhempaan XML-standardiin \cite{nurseitov_comparison_nodate}. Uusi raportointityökalu hyödyntää lähdedatan hankkimisessa Sovelia PLM-järjestelmän olemassa olevaa hakumoottoria, joka palauttaa itsessään JSON-muotoista dataa, joten muuntamisvaihetta erilliseen dataformaattiin ei tarvita.

% Kuva raporteista? Esim. Scubalight BOM PDF ja Excel?

Sovelia PLM:n uusi raportointityökalu tuottaa raportteja PDF-, Excel ja ZIP-muodoissa eli se tukee tarkasteltujen raportointityökalujen suosituimpia formaatteja. PDF-raportit tarjoavat taulukoidun tuotedatan lisäksi myös piirakkakuvaajia tuoterakenteen eri tietokenttien arvojen jakaumasta, Excel-raportit mahdollistavat rakenteellisen esittämisen tasoittain ja tarjoavat laskentamahdollisuuksia. ZIP-raportit ovat taas kokoavat tuoterakenteesta useita PDF- tai Excel-raportteja ja lisäävät arkistoon myös tuoteobjektien liitetiedostoja, kuten piirustuksia. Uusi raportointimoottori tukee myös HTML-raportteja, mutta ne eivät kuulu ensimmäisen version vaatimusmäärittelyyn, eivätkä siten ole osa ensimmäistä julkaistavaa versiota.

% Kuva käyttöliittymästä?

Raporttipohjien konfigurointi tapahtuu uudella raportointityökalulla tapahtuu Sovelia PLM:n web-käyttöliittymässä pääkäyttäjän toimesta. Käyttöliittymässä voidaan luoda, muokata ja poistaa raportteja, sekä säätää interaktiivisten valikoiden avulla esimerkiksi tietokentät, joista muodostetaan raportin kuvaajat ja taulukot. Raportointimoottori osaa myös laskea lisäkenttiä, kuten esimerkiksi osien kokonaissummia tuoterakenteessa. Tuoterakenteista voidaan myös suodattaa pois määritellyn ominaisuuden mukaisesti tuoteobjekteja, tai raportissa esitettävä tuoterakenne voidaan asettaa koostumaan vain tietyn ominaisuuden omaavista tuoteobjekteista. ZIP-raporttipaketteihin voidaan valita halutut tuoteobjektien liitetiedostojen tiedostomuodot. Myös PDF-raportteihin on mahdollista liittää tuotteiden ja osien piirustuksia PDF-muodossa. ZIP-paketin "aliraportiksi" voidaan valita mikä tahansa järjestelmään konfiguroitu Excel- tai PDF-raportti. Aliraportilla tarkoitetaan tässä raporttia, joka ajetaan käyttöliittymässä annettujen kriteerien mukaisesti tuoterakenteen eri alirakenteille tai osille.

% Kuva ZIP-raportin rakenteesta?

\subsubsection{Uuden raportointityökalun tulevaisuus}

Uusi raportointimoottori Node.js -palvelimella on osoittautunut testausvaiheessa tehokkaaksi ja nopeaksi keinoksi muodostaa raportteja. Suoritusajat ovat huomattavasti vanhaa raportointityökalua nopeampia ja muistinkäyttö palvelimella on vähäisempää. Konfigurointia on pyritty tekemään loppukäyttäjille helpommaksi tarjoamalla graafinen web-käyttöliittymä raporttien konfigurointiin ja raportteja voidaan luoda suoraan Sovelia PLM:n valikoista. Tulevaisuuden ominaisuuksina on tarkoitus kehittää graafinen suunnittelutyökalu PDF-raporteille, lähettää monitorointidataa raportointimoottorin suorituskyvystä sekä tarjota loppukäyttäjille enemmän mahdollisuuksia vaikuttaa raporttien ulkoasuun.



%\input{file_name_of_chapter_x}
%\input{file_name_of_chapter_y}

% The thesis main content ends here.

\printbibliography

\begin{comment}
Important! Create the appendix chapters with command \textbackslash appchapter\{some
name\} instead of \textbackslash chapter\{some name\} for the automagic
page counting to work!
\end{comment}


\appchapter{Yleistä raportointityökaluista}

\begin{landscape}
\begin{table}
    \centering
    \scriptsize  % Adjust font size to scriptsize
    \resizebox{\linewidth}{!}{%
        \begin{tabular}{|p{2cm}|p{2cm}|p{2.5cm}|p{2cm}|p{3cm}|p{2cm}|}
            \hline
            \bfseries Raportointi-työkalu & \bfseries Lisenssi & \bfseries Datalähde & \bfseries Ohjelmointi-kieli & \bfseries Raporttipohjien konfigurointi & \bfseries Asennus \\ \hline
            ActiveReportsJS & kaupallinen & JSON & JavaScript & UI /JSON & paikallinen \\ \hline
            SAP Business Objects & kaupallinen & XML ja JSON & Java & UI & paikallinen \\ \hline
            SAP Crystal Reports & kaupallinen & tiedostot mm. XLSX, TXT, XML ja tietokannat (SQL) & Java & UI & paikallinen \\ \hline
            Oracle BI Publisher & kaupallinen & SQL, XML, XLSX ja Oracle BI:n omat formaatit & Java & UI ja raportointipohjat & paikallinen \\ \hline
            JasperReports & avoin lähdekoodi, mahdollisuus kaupalliseen lisenssiin & JavaBeans, Java object, XML ja mukautetut tietolähteet & Java & UI / XML tai .jasper & paikallinen \\ \hline
            BIRT Project & avoin lähdekoodi & XML, TXT ja tietokannat (SQL) & Java & UI / XML & paikallinen \\
            \hline
        \end{tabular}%
    }
    \captionsetup{skip=6pt}  % Adjust caption spacing
    \caption{Yleistä raportointityökaluista. \textit{Lisää lähteet!}\label{tab:Yleista}}
\end{table}
\end{landscape}

\appchapter{Raportointityökalujen ulostuloformaatit}

\begin{landscape}
\begin{table}
    \centering
    \scriptsize  % Adjust font size to scriptsize
    \resizebox{\linewidth}{!}{%
    \begin{tabular}{|l|l|l|l|l|l|l|}
        \hline
        \bfseries Raportointityökalu & \bfseries PDF (.pdf) & \bfseries XLSX (.xlsx, .xls) & \bfseries HTML (.html) & \bfseries Word (.docx, .rtf) & \bfseries RPT (.rpt) \\ \hline
        ActiveReportsJS & x & & x & & \\ \hline
        SAP BusinessObjects & x & x & & x & x \\ \hline
        SAP Crystal Reports & x & x & x & x & x \\ \hline
        Oracle BI Publisher & x & x & & x & x \\ \hline
        JasperReports & x & x & x & x & \\ \hline
        BIRT Project & x & x & x & x & \\ \hline
    \end{tabular}%
    }
    \captionsetup{skip=6pt}  % Adjust caption spacing
    \caption{Raportointityökalujen ulostuloformaatit. \textit{Lisää lähteet!}\label{tab:Formaatit}}
\end{table}
\end{landscape}

\end{document}
